\section{Conclusion}
\label{sec:conclusion}

\paragraph{Summary.}
We studied the question of how PPO should interact with the multi-step
stochastic process inside a flow-based generative policy.  We formalized
per-step PPO, which applies clipped importance-sampling ratios independently at
each generative step, and introduced weighted credit-assignment schemes---uniform,
learned-global, and state-dependent---to modulate gradient signal across steps.
\todo{State the main empirical finding in 1--2 sentences.  E.g., ``Per-step PPO
with learned global weights consistently outperformed holistic PPO on
HalfCheetah-v5 and Hopper-v5, with the largest gains observed for $K{=}8$
flow steps.''}

\paragraph{Limitations.}
Our study has several limitations:
\begin{itemize}[nosep,leftmargin=*]
  \item \textbf{Computational overhead.}  Per-step PPO requires storing
        intermediate latents $\{\mathbf{z}_{t,k}\}$ for every transition and
        computing $K$ ratios per sample, increasing memory and compute relative
        to holistic PPO.
  \item \textbf{Limited environments.}  We evaluated on two MuJoCo locomotion
        tasks.  Results may differ on higher-dimensional or sparse-reward tasks
        such as AntMaze or dexterous manipulation.
  \item \textbf{Fixed noise schedule.}  We used a constant $\sigma$ across all
        steps; an adaptive or learned noise schedule could interact differently
        with per-step clipping.
  \item \textbf{No theoretical guarantees.}  We lack a formal analysis of when
        per-step clipping provides better trust-region properties than holistic
        clipping.
\end{itemize}

\paragraph{Future Work.}
Several directions are promising:
\begin{itemize}[nosep,leftmargin=*]
  \item \textbf{Learned noise schedule.}  Jointly learning $\sigma_k$ per step
        could improve both policy expressiveness and the quality of per-step
        ratios.
  \item \textbf{Scaling to larger $K$.}  With many flow steps ($K \geq 16$),
        the advantage of per-step clipping should grow; evaluating this regime
        is important.
  \item \textbf{Broader benchmarks.}  Testing on image-based manipulation,
        multi-task settings, and sparse-reward environments would strengthen
        generality claims.
  \item \textbf{Theoretical analysis.}  Characterizing the bias-variance
        trade-off of per-step vs.\ holistic clipping, and conditions under
        which weighted credit assignment provably improves convergence, would
        complement our empirical results.
\end{itemize}
